
\begin{thebibliography}{9}

\bibitem{kobayashi2024lowrank}
Kobayashi, S., Akram, Y., \& von Oswald, J. (2024). Weight decay induces low-rank attention layers. \textit{Advances in Neural Information Processing Systems (NeurIPS)}. \texttt{arXiv:2410.23819}.

\bibitem{srebroshraibman2005}
Srebro, N., \& Shraibman, A. (2005). Rank, trace-norm and max-norm. \textit{Conference on Learning Theory (COLT)}, 545-560.

\bibitem{rechtfazelparrilo2010}
Recht, B., Fazel, M., \& Parrilo, P. A. (2010). Guaranteed minimum-rank solutions of linear matrix equations via nuclear norm minimization. \textit{SIAM Review}, 52(3), 471-501.

\bibitem{tu2016procrustes}
Tu, S., Boumal, N., Bhojanapalli, S., Sanghavi, S., \& Recht, B. (2016). Low-rank solutions of linear matrix equations via Procrustes flow. \textit{International Conference on Machine Learning (ICML)}.

\bibitem{burermonteiro2003}
Burer, S., \& Monteiro, R. D. (2003). A nonlinear programming algorithm for solving semidefinite programs via low-rank factorization. \textit{Mathematical Programming}, 95(2), 329-357.
\bibitem{bm2003}
  Burer, S., \& Monteiro, R.\,D.\,C. (2003).
  A nonlinear programming algorithm for solving semidefinite programs via low-rank
  factorization.
  \emph{Mathematical Programming}, 95(2), 329--357.

\bibitem{bm2005}
  Burer, S., \& Monteiro, R.\,D.\,C. (2005).
  Local minima and convergence in low-rank semidefinite programming.
  \emph{Mathematical Programming}, 103(3), 427--444.

\bibitem{boumal2016}
  Boumal, N., Voroninski, V., \& Bandeira, A.\,S. (2016).
  The non-convex Burer-Monteiro approach works on smooth semidefinite programs.
  \emph{Advances in Neural Information Processing Systems}, 29.

\bibitem{gunasekar2017implicit}
  Gunasekar, S., Woodworth, B., Bhojanapalli, S., Neyshabur, B., \& Srebro, N. (2017).
  Implicit regularization in matrix factorization.
  \emph{Advances in Neural Information Processing Systems}, 30.

\bibitem{iclr2023}
  Zhu, M., Hayashi, H., \& Ma, H. (2023).
  Scaling convex neural networks with Burer-Monteiro factorization.
  \emph{International Conference on Learning Representations}.

\bibitem{lee2016sgd}
  Lee, J.\,D., Simchowitz, M., Jordan, M.\,I., \& Recht, B. (2016).
  Gradient descent only converges to minimizers.
  \emph{Conference on Learning Theory}, 1246--1257.

\bibitem{candes2009}
  Cand\`{e}s, E.\,J., \& Recht, B. (2009).
  Exact matrix completion via convex optimization.
  \emph{Foundations of Computational Mathematics}, 9(6), 717--772.

\bibitem{vaswani2017}
  Vaswani, A., Shazeer, N., Parmar, N., Uszkoreit, J., Jones, L., Gomez, A.\,N.,
  Kaiser, L., \& Polosukhin, I. (2017).
  Attention is all you need.
  \emph{Advances in Neural Information Processing Systems}, 30.

\bibitem{srebro2005}
  Srebro, N., \& Shraibman, A. (2005).
  Rank, trace-norm and max-norm.
  \emph{Conference on Learning Theory}, 545--560.

\bibitem{kobayashi2024}
  Kobayashi, S., Akram, Y., \& von Oswald, J. (2024).
  Weight decay induces low-rank attention layers.
  \emph{Advances in Neural Information Processing Systems}.

\bibitem{arora2019}
  Arora, S., Cohen, N., Hu, W., \& Luo, Y. (2019).
  Implicit regularization in deep matrix factorization.
  \emph{Advances in Neural Information Processing Systems}, 32.

\bibitem{barvinok1995}
  Barvinok, A. (1995).
  Problems of distance geometry and convex properties of quadratic maps.
  \emph{Discrete \& Computational Geometry}, 13, 189--202.

\bibitem{pataki1998}
  Pataki, G. (1998).
  On the rank of extreme matrices in semidefinite programs and the multiplicity of
  optimal eigenvalues.
  \emph{Mathematics of Operations Research}, 23(2), 339--358.

\bibitem{waldspurger2020}
  Waldspurger, I., \& Waters, A. (2020).
  Rank optimality for the Burer-Monteiro factorization.
  \emph{SIAM Journal on Optimization}, 30(3), 2577--2602.
\end{thebibliography}

